% =============================================================================
% PoCA Vertex Reconstruction — Mathematical Derivation & Verification Guide
% Include in parent document with: % =============================================================================
% PoCA Vertex Reconstruction — Mathematical Derivation & Verification Guide
% Include in parent document with: % =============================================================================
% PoCA Vertex Reconstruction — Mathematical Derivation & Verification Guide
% Include in parent document with: % =============================================================================
% PoCA Vertex Reconstruction — Mathematical Derivation & Verification Guide
% Include in parent document with: \input{poca_vertex_reconstruction}
% Required packages in parent preamble:
%   \usepackage{amsmath, amssymb, booktabs, array}
% =============================================================================

% =============================================================================
\subsection{Track Parameterisation}
% =============================================================================

Each pion track is a straight line in 3D, fitted by two independent
linear regressions: one in the $xz$-plane and one in the $yz$-plane.

\begin{equation}
    x(z) = b_x + a_x z, \qquad y(z) = b_y + a_y z
\end{equation}

\noindent
Here $a_x = dx/dz$ and $a_y = dy/dz$ are the fitted track slopes (the rate of transverse displacement per unit distance along the beam axis), and $b_x$, $b_y$ are the track intercepts --- the transverse positions of the track at $z = 0$.

The unit direction vector and reference point (evaluated at $z = 0$, used purely as a mathematical anchor for the fit --- not a physical location) are:

\begin{equation}
    \vec{d} = \bigl(a_x,\; a_y,\; 1\bigr)^{\!\text{unit}}, \qquad
    \vec{o} = (b_x,\; b_y,\; 0)
\end{equation}

Any point on the track is then $\vec{r}(s) = \vec{o} + s\,\vec{d}$, where $s$ is a scalar arc-length parameter: stepping $s$ forward or backward moves along the track in units of centimetres (since $\vec{d}$ is a unit vector).

\subsubsection{Weighted Least-Squares for \texorpdfstring{$a_x, b_x$}{ax, bx}}

To compute the fitted slopes and intercepts by hand from the hit positions,
use weighted least squares. For $N$ hit points $(z_i, x_i)$ --- where $N$ is the total number of hits contributing to the fit, $z_i$ is the beam-axis position of hit $i$, and $x_i$ is the measured transverse coordinate --- with weights $w_i = 1/\sigma_i^2$, where $\sigma_i$ is the positional resolution of the detector element that produced hit $i$ (values given in the table below):

\begin{equation}
    a_x = \frac{\langle zx \rangle - \langle z \rangle\langle x \rangle}{\langle z^2 \rangle - \langle z \rangle^2},
    \qquad
    b_x = \langle x \rangle - a_x\,\langle z \rangle,
    \qquad \text{where} \quad
    \langle f \rangle = \frac{\sum_i w_i f_i}{\sum_i w_i}
\end{equation}

The same procedure gives $a_y, b_y$ from the $(z_i, y_i)$ points.

\subsubsection{Detector Measurement Uncertainties}

The per-hit uncertainties $\sigma_i$ used as weights in the track fit are:

\begin{center}
\begin{tabular}{llll}
    \toprule
    \textbf{Detector} & \textbf{Device IDs} & \textbf{Axis measured} & \textbf{$\sigma$ (cm)} \\
    \midrule
    T1 (straws along $Y$)        & 1--488      & $X$ only            & 0.40 \\
    T2 (straws along $X$)        & 489--976    & $Y$ only            & 0.40 \\
    T3 (stereo $+45°$)           & 977--1464   & $X$ and $Y$         & $0.40\sqrt{2} \approx 0.566$ \\
    T4 (stereo $-45°$)           & 1465--1952  & $X$ and $Y$         & $0.40\sqrt{2} \approx 0.566$ \\
    FRI Wall 1 (inner strips)    & 2001--2010, 2017--2026 & $X$ only & 1.50 \\
    FRI Wall 1 (outer strips)    & 2011--2016  & $X$ only            & 3.00 \\
    FRI Wall 2 (inner strips)    & 2027--2036, 2043--2052 & $Y$ only & 1.50 \\
    FRI Wall 2 (outer strips)    & 2037--2042  & $Y$ only            & 3.00 \\
    TOF wall bars                & 2053--2070  & \textit{excluded from fit} & --- \\
    \bottomrule
\end{tabular}
\end{center}

\noindent
The TOF bars are excluded from the track fit because their large positional
uncertainties ($\pm 6$ cm in $X$, $\pm 10$ cm in $Y$) at the $\sim\!460$ cm
lever arm would bias the fitted slope and degrade backward-extrapolated vertex
resolution. TOF hits are used only for the velocity measurement.

% =============================================================================
\subsection{Point of Closest Approach (PoCA) --- Full Derivation}
% =============================================================================

Given two fitted straight-line tracks:

\begin{equation}
    \vec{r}_1(s) = \vec{p}_1 + s\,\vec{v}_1, \qquad
    \vec{r}_2(t) = \vec{p}_2 + t\,\vec{v}_2,
\end{equation}

\noindent
where $\vec{p}_{1,2}$ are the mathematical reference points on each track, defined as the position where the fitted line passes through the $z = 0$ plane (this is not the decay vertex or the kaon production point --- it is simply the intercept of the extrapolated track line, used as a convenient algebraic anchor), $\vec{v}_{1,2}$ are the unit direction vectors for the $\pi^+$ and $\pi^-$ tracks respectively (corresponding to $\vec{d}$ above), and $s$, $t$ are free scalar arc-length parameters that index positions along each track. The goal is to find the values of $s$ and $t$ that minimise the distance $\bigl|\vec{r}_1(s) - \vec{r}_2(t)\bigr|$.

\subsubsection{Setup}

Define:

\begin{equation}
    \vec{w}_0 = \vec{p}_1 - \vec{p}_2
\end{equation}

\noindent
The vector $\vec{w}_0$ is the displacement from the $\pi^-$ reference point to the $\pi^+$ reference point; it encodes the initial separation between the two track lines at $z = 0$.

Define also the five scalar dot products:

\begin{align}
    a &= \vec{v}_1 \cdot \vec{v}_1, &
    b &= \vec{v}_1 \cdot \vec{v}_2, &
    c &= \vec{v}_2 \cdot \vec{v}_2, \\[4pt]
    d &= \vec{v}_1 \cdot \vec{w}_0, &
    e &= \vec{v}_2 \cdot \vec{w}_0. & &
\end{align}

\noindent
Since $\vec{v}_{1,2}$ are unit vectors, $a = c = 1$ always. The quantity $b = \cos\theta$ is the cosine of the opening angle $\theta$ between the two pion tracks. The scalars $d$ and $e$ project the inter-track separation vector $\vec{w}_0$ onto the $\pi^+$ and $\pi^-$ direction vectors respectively; they determine how the gap between the reference points must be closed by stepping along each track.

\subsubsection{Minimisation}

Setting $\partial/\partial s$ and $\partial/\partial t$ of
$|\vec{r}_1(s) - \vec{r}_2(t)|^2$ to zero gives the $2 \times 2$ linear system:

\begin{equation}
    \begin{pmatrix} a & -b \\ b & -c \end{pmatrix}
    \begin{pmatrix} s_c \\ t_c \end{pmatrix}
    = \begin{pmatrix} -d \\ -e \end{pmatrix}
\end{equation}

Solving (Cramer's rule), with denominator $D = ac - b^2 = 1 - \cos^2\theta = \sin^2\theta$:

\begin{equation}
    \boxed{s_c = \frac{be - cd}{D}, \qquad t_c = \frac{ae - bd}{D}}
\end{equation}

\noindent
The denominator $D = \sin^2\theta$ is a measure of how non-parallel the two tracks are; it reaches its maximum of 1 for perpendicular tracks and vanishes as the opening angle $\theta \to 0$. The solved parameters $s_c$ and $t_c$ are the arc-length values at which each track is closest to the other --- stepping $s_c$ along track 1 and $t_c$ along track 2 brings both tracks to their point of minimum mutual separation. If $|D| < 10^{-14}$ the lines are parallel and no unique solution exists.

\subsubsection{Closest Points and Decay Vertex}

The closest point on each track is:

\begin{equation}
    \vec{q}_1 = \vec{p}_1 + s_c\,\vec{v}_1, \qquad
    \vec{q}_2 = \vec{p}_2 + t_c\,\vec{v}_2
\end{equation}

\noindent
The points $\vec{q}_1$ and $\vec{q}_2$ are the feet of the common perpendicular between the two track lines --- the positions on the $\pi^+$ and $\pi^-$ tracks respectively that are nearest to each other. In a perfect detector with infinitely precise tracking, both pions would originate from the same point and $\vec{q}_1 = \vec{q}_2$ exactly.

The reconstructed decay vertex is the midpoint of the shortest connecting
segment:

\begin{equation}
    \boxed{\vec{V}_\text{reco} = \tfrac{1}{2}\bigl(\vec{q}_1 + \vec{q}_2\bigr)}
\end{equation}

The \emph{PoCA separation} (a quality indicator) is:

\begin{equation}
    \Delta = |\vec{q}_1 - \vec{q}_2|
\end{equation}

For a clean decay event occurring at a well-defined vertex, $\Delta$ should be
close to zero. Large $\Delta$ indicates nearly-parallel tracks, a track mismatch,
or a poor fit.

% =============================================================================
\subsection{Physical Interpretation of \texorpdfstring{$s_c$}{sc} and \texorpdfstring{$t_c$}{tc}}
% =============================================================================

Because $\vec{v}_{1,2}$ are unit vectors ($|\vec{v}| = 1$), the parameters
$s_c$ and $t_c$ are the \textbf{arc-lengths along each track} from the reference
point $\vec{o}$ to the PoCA point. A negative value means the vertex lies
upstream of $z = 0$.

The $z$-coordinate of the PoCA point on track 1 is:

\begin{equation}
    z_{\text{PoCA},1} = o_z + s_c \cdot d_z
\end{equation}

\noindent
where $o_z = 0$ is the $z$-component of the reference point $\vec{o}$ (the mathematical intercept anchor of the fit, not the decay vertex location), and $d_z$ is the $z$-component of the unit direction vector $\vec{d}$. For a near-forward track $d_z \approx 1$, so $s_c \approx z_{\text{PoCA},1}$ numerically.
The corresponding transverse position is:

\begin{equation}
    x_{\text{PoCA},1} = o_x + s_c \cdot d_x, \qquad
    y_{\text{PoCA},1} = o_y + s_c \cdot d_y
\end{equation}

\noindent
where $o_x = b_x$ and $o_y = b_y$ are the fitted intercepts, and $d_x$, $d_y$ are the $x$ and $y$ components of the unit direction vector.

% =============================================================================
\subsection{Selecting a Clean Event for Manual Verification}
% =============================================================================

Enable \texttt{DIAG\_RECO\_MODE = true} in \texttt{KLong\_save\_vectors.C} and
scan the printed output for an event satisfying all of the following criteria:

\begin{center}
\begin{tabular}{lll}
    \toprule
    \textbf{Diagnostic} & \textbf{Item in output} & \textbf{Target} \\
    \midrule
    PoCA separation $\Delta$          & \texttt{[1]} & $< 1$ cm \\
    Back-propagated decay time spread & \texttt{[2]} & $|\text{diff}| < 0.1$ ns \\
    Vertex $z$ residual $\delta z$    & \texttt{[3]} & $|\delta z| < 5$ cm \\
    Track slope errors                & \texttt{[4]} & $|\Delta a_x|, |\Delta a_y| < 0.01$ \\
    Tracker coverage                  & hit counts   & T1, T2, T3, T4 all $> 0$ \\
    \bottomrule
\end{tabular}
\end{center}

From a clean event, record the following from the \texttt{DIAG} printout:
\begin{itemize}
    \item All raw hit positions per detector plane (from the \texttt{DEBUG} block).
    \item \texttt{orig} $(b_x, b_y, 0)$ and \texttt{dir} $(d_x, d_y, d_z)$ for each fitted track.
    \item The printed \texttt{reco\_vertex} and \texttt{true\_vertex}.
\end{itemize}

% =============================================================================
\subsection{Manual Calculation Walkthrough}
% =============================================================================

Given the track parameters extracted from a DIAG event for the $\pi^+$ and
$\pi^-$ tracks:

\begin{align}
    \vec{p}_{1} &= (b_x^{\pi^+},\; b_y^{\pi^+},\; 0), &
    \vec{v}_{1} &= (a_x^{\pi^+},\; a_y^{\pi^+},\; 1)^{\text{unit}} \\
    \vec{p}_{2} &= (b_x^{\pi^-},\; b_y^{\pi^-},\; 0), &
    \vec{v}_{2} &= (a_x^{\pi^-},\; a_y^{\pi^-},\; 1)^{\text{unit}}
\end{align}

Proceed step by step:

\begin{enumerate}
    \item Compute $\vec{w}_0 = \vec{p}_1 - \vec{p}_2$.

    \item Since $\vec{v}_{1,2}$ are unit vectors: $a = c = 1$, so
    \begin{equation}
        D = 1 - b^2 = \sin^2\theta,
    \end{equation}
    where $\theta$ is the opening angle between the $\pi^+$ and $\pi^-$ tracks --- the angle between their fitted direction vectors in 3D --- and $b = \cos\theta = \vec{v}_1 \cdot \vec{v}_2$ is its cosine, computed directly from the track fit outputs.

    \item Compute:
    \begin{equation}
        d = \vec{v}_1 \cdot \vec{w}_0, \qquad e = \vec{v}_2 \cdot \vec{w}_0
    \end{equation}

    \item Solve for the PoCA parameters:
    \begin{equation}
        s_c = \frac{be - d}{D}, \qquad t_c = \frac{e - bd}{D}
    \end{equation}

    \item Compute the two closest points:
    \begin{equation}
        \vec{q}_1 = \vec{p}_1 + s_c\,\vec{v}_1, \qquad
        \vec{q}_2 = \vec{p}_2 + t_c\,\vec{v}_2
    \end{equation}

    \item The decay vertex estimate is:
    \begin{equation}
        \vec{V}_\text{reco} = \tfrac{1}{2}\bigl(\vec{q}_1 + \vec{q}_2\bigr)
    \end{equation}

    This should match the \texttt{reco\_vertex} in the DIAG output to three
    decimal places.
\end{enumerate}

% =============================================================================
\subsection{Why \texorpdfstring{$z$}{z} is the Hard Coordinate}
% =============================================================================

The transverse coordinates $x$ and $y$ at the vertex are
\textbf{directly constrained} by the T1 and T2 measurements respectively.
The vertex $z$-position, however, is \textbf{never directly measured} --- it
is determined entirely by the fitted slopes $a_x, a_y$ and the PoCA algebra.

The $z$ error propagates approximately as:

\begin{equation}
    \delta z_\text{vtx} \;\sim\; \frac{\sigma_\text{slope}}{\Delta a^\text{opening}}
\end{equation}

where $\Delta a^\text{opening}$ is the difference in slopes between the
$\pi^+$ and $\pi^-$ tracks and $\sigma_\text{slope}$ is the typical
uncertainty on the fitted slope.

When the opening angle is small (high-momentum kaon, both pions nearly
collinear), $\Delta a^\text{opening}$ is tiny and the $z$ uncertainty diverges.
This is precisely why the vertex summary shows:

\begin{center}
\begin{tabular}{lcc}
    \toprule
    \textbf{Coordinate} & $\mu$ (cm) & $\sigma$ (cm) \\
    \midrule
    $\delta x$ & $\sim -0.05$ & $\sim 0.4$ \\
    $\delta y$ & $\sim 0.00$  & $\sim 0.35$ \\
    $\delta z$ & $\sim +1$--5 & $\sim 7$--12 \\
    \bottomrule
\end{tabular}
\end{center}

\noindent The $z$ resolution is 20--30$\times$ worse than the transverse
resolution, and the positive $\mu_z$ bias (downstream shift) indicates
systematic overestimation of the vertex $z$ position, which propagates
directly into an underestimation of the reconstructed kaon momentum.

% =============================================================================
\subsection{Quick-Start Verification Recipe}
% =============================================================================

The most direct path to a single verified event is:

\begin{enumerate}
    \item In \texttt{KLong\_save\_vectors.C}, set:
    \begin{verbatim}
const bool DEBUG_MODE       = true;
const int  DEBUG_MAX_EVENTS = 1;
const bool DIAG_RECO_MODE   = true;
const int  DIAG_MAX_EVENTS  = 1;
    \end{verbatim}

    \item Run on one simulation file. The first successfully reconstructed
    event will print the full hit list, fitted track parameters, and the
    5-item diagnostic checklist.

    \item Extract \texttt{orig} and \texttt{dir} for each pion track from
    the output.

    \item Work through the 6-step PoCA calculation in Section~5.

    \item Compare your computed $\vec{V}_\text{reco}$ to the printed
    \texttt{reco\_vertex} --- they should agree to $\lesssim 10^{-3}$ cm.

    \item Compare $\vec{V}_\text{reco}$ to the printed \texttt{true\_vertex}
    (from Geant4 MC truth) to quantify the reconstruction resolution for
    that event.
\end{enumerate}

% Required packages in parent preamble:
%   \usepackage{amsmath, amssymb, booktabs, array}
% =============================================================================

% =============================================================================
\subsection{Track Parameterisation}
% =============================================================================

Each pion track is a straight line in 3D, fitted by two independent
linear regressions: one in the $xz$-plane and one in the $yz$-plane.

\begin{equation}
    x(z) = b_x + a_x z, \qquad y(z) = b_y + a_y z
\end{equation}

\noindent
Here $a_x = dx/dz$ and $a_y = dy/dz$ are the fitted track slopes (the rate of transverse displacement per unit distance along the beam axis), and $b_x$, $b_y$ are the track intercepts --- the transverse positions of the track at $z = 0$.

The unit direction vector and reference point (evaluated at $z = 0$, used purely as a mathematical anchor for the fit --- not a physical location) are:

\begin{equation}
    \vec{d} = \bigl(a_x,\; a_y,\; 1\bigr)^{\!\text{unit}}, \qquad
    \vec{o} = (b_x,\; b_y,\; 0)
\end{equation}

Any point on the track is then $\vec{r}(s) = \vec{o} + s\,\vec{d}$, where $s$ is a scalar arc-length parameter: stepping $s$ forward or backward moves along the track in units of centimetres (since $\vec{d}$ is a unit vector).

\subsubsection{Weighted Least-Squares for \texorpdfstring{$a_x, b_x$}{ax, bx}}

To compute the fitted slopes and intercepts by hand from the hit positions,
use weighted least squares. For $N$ hit points $(z_i, x_i)$ --- where $N$ is the total number of hits contributing to the fit, $z_i$ is the beam-axis position of hit $i$, and $x_i$ is the measured transverse coordinate --- with weights $w_i = 1/\sigma_i^2$, where $\sigma_i$ is the positional resolution of the detector element that produced hit $i$ (values given in the table below):

\begin{equation}
    a_x = \frac{\langle zx \rangle - \langle z \rangle\langle x \rangle}{\langle z^2 \rangle - \langle z \rangle^2},
    \qquad
    b_x = \langle x \rangle - a_x\,\langle z \rangle,
    \qquad \text{where} \quad
    \langle f \rangle = \frac{\sum_i w_i f_i}{\sum_i w_i}
\end{equation}

The same procedure gives $a_y, b_y$ from the $(z_i, y_i)$ points.

\subsubsection{Detector Measurement Uncertainties}

The per-hit uncertainties $\sigma_i$ used as weights in the track fit are:

\begin{center}
\begin{tabular}{llll}
    \toprule
    \textbf{Detector} & \textbf{Device IDs} & \textbf{Axis measured} & \textbf{$\sigma$ (cm)} \\
    \midrule
    T1 (straws along $Y$)        & 1--488      & $X$ only            & 0.40 \\
    T2 (straws along $X$)        & 489--976    & $Y$ only            & 0.40 \\
    T3 (stereo $+45°$)           & 977--1464   & $X$ and $Y$         & $0.40\sqrt{2} \approx 0.566$ \\
    T4 (stereo $-45°$)           & 1465--1952  & $X$ and $Y$         & $0.40\sqrt{2} \approx 0.566$ \\
    FRI Wall 1 (inner strips)    & 2001--2010, 2017--2026 & $X$ only & 1.50 \\
    FRI Wall 1 (outer strips)    & 2011--2016  & $X$ only            & 3.00 \\
    FRI Wall 2 (inner strips)    & 2027--2036, 2043--2052 & $Y$ only & 1.50 \\
    FRI Wall 2 (outer strips)    & 2037--2042  & $Y$ only            & 3.00 \\
    TOF wall bars                & 2053--2070  & \textit{excluded from fit} & --- \\
    \bottomrule
\end{tabular}
\end{center}

\noindent
The TOF bars are excluded from the track fit because their large positional
uncertainties ($\pm 6$ cm in $X$, $\pm 10$ cm in $Y$) at the $\sim\!460$ cm
lever arm would bias the fitted slope and degrade backward-extrapolated vertex
resolution. TOF hits are used only for the velocity measurement.

% =============================================================================
\subsection{Point of Closest Approach (PoCA) --- Full Derivation}
% =============================================================================

Given two fitted straight-line tracks:

\begin{equation}
    \vec{r}_1(s) = \vec{p}_1 + s\,\vec{v}_1, \qquad
    \vec{r}_2(t) = \vec{p}_2 + t\,\vec{v}_2,
\end{equation}

\noindent
where $\vec{p}_{1,2}$ are the mathematical reference points on each track, defined as the position where the fitted line passes through the $z = 0$ plane (this is not the decay vertex or the kaon production point --- it is simply the intercept of the extrapolated track line, used as a convenient algebraic anchor), $\vec{v}_{1,2}$ are the unit direction vectors for the $\pi^+$ and $\pi^-$ tracks respectively (corresponding to $\vec{d}$ above), and $s$, $t$ are free scalar arc-length parameters that index positions along each track. The goal is to find the values of $s$ and $t$ that minimise the distance $\bigl|\vec{r}_1(s) - \vec{r}_2(t)\bigr|$.

\subsubsection{Setup}

Define:

\begin{equation}
    \vec{w}_0 = \vec{p}_1 - \vec{p}_2
\end{equation}

\noindent
The vector $\vec{w}_0$ is the displacement from the $\pi^-$ reference point to the $\pi^+$ reference point; it encodes the initial separation between the two track lines at $z = 0$.

Define also the five scalar dot products:

\begin{align}
    a &= \vec{v}_1 \cdot \vec{v}_1, &
    b &= \vec{v}_1 \cdot \vec{v}_2, &
    c &= \vec{v}_2 \cdot \vec{v}_2, \\[4pt]
    d &= \vec{v}_1 \cdot \vec{w}_0, &
    e &= \vec{v}_2 \cdot \vec{w}_0. & &
\end{align}

\noindent
Since $\vec{v}_{1,2}$ are unit vectors, $a = c = 1$ always. The quantity $b = \cos\theta$ is the cosine of the opening angle $\theta$ between the two pion tracks. The scalars $d$ and $e$ project the inter-track separation vector $\vec{w}_0$ onto the $\pi^+$ and $\pi^-$ direction vectors respectively; they determine how the gap between the reference points must be closed by stepping along each track.

\subsubsection{Minimisation}

Setting $\partial/\partial s$ and $\partial/\partial t$ of
$|\vec{r}_1(s) - \vec{r}_2(t)|^2$ to zero gives the $2 \times 2$ linear system:

\begin{equation}
    \begin{pmatrix} a & -b \\ b & -c \end{pmatrix}
    \begin{pmatrix} s_c \\ t_c \end{pmatrix}
    = \begin{pmatrix} -d \\ -e \end{pmatrix}
\end{equation}

Solving (Cramer's rule), with denominator $D = ac - b^2 = 1 - \cos^2\theta = \sin^2\theta$:

\begin{equation}
    \boxed{s_c = \frac{be - cd}{D}, \qquad t_c = \frac{ae - bd}{D}}
\end{equation}

\noindent
The denominator $D = \sin^2\theta$ is a measure of how non-parallel the two tracks are; it reaches its maximum of 1 for perpendicular tracks and vanishes as the opening angle $\theta \to 0$. The solved parameters $s_c$ and $t_c$ are the arc-length values at which each track is closest to the other --- stepping $s_c$ along track 1 and $t_c$ along track 2 brings both tracks to their point of minimum mutual separation. If $|D| < 10^{-14}$ the lines are parallel and no unique solution exists.

\subsubsection{Closest Points and Decay Vertex}

The closest point on each track is:

\begin{equation}
    \vec{q}_1 = \vec{p}_1 + s_c\,\vec{v}_1, \qquad
    \vec{q}_2 = \vec{p}_2 + t_c\,\vec{v}_2
\end{equation}

\noindent
The points $\vec{q}_1$ and $\vec{q}_2$ are the feet of the common perpendicular between the two track lines --- the positions on the $\pi^+$ and $\pi^-$ tracks respectively that are nearest to each other. In a perfect detector with infinitely precise tracking, both pions would originate from the same point and $\vec{q}_1 = \vec{q}_2$ exactly.

The reconstructed decay vertex is the midpoint of the shortest connecting
segment:

\begin{equation}
    \boxed{\vec{V}_\text{reco} = \tfrac{1}{2}\bigl(\vec{q}_1 + \vec{q}_2\bigr)}
\end{equation}

The \emph{PoCA separation} (a quality indicator) is:

\begin{equation}
    \Delta = |\vec{q}_1 - \vec{q}_2|
\end{equation}

For a clean decay event occurring at a well-defined vertex, $\Delta$ should be
close to zero. Large $\Delta$ indicates nearly-parallel tracks, a track mismatch,
or a poor fit.

% =============================================================================
\subsection{Physical Interpretation of \texorpdfstring{$s_c$}{sc} and \texorpdfstring{$t_c$}{tc}}
% =============================================================================

Because $\vec{v}_{1,2}$ are unit vectors ($|\vec{v}| = 1$), the parameters
$s_c$ and $t_c$ are the \textbf{arc-lengths along each track} from the reference
point $\vec{o}$ to the PoCA point. A negative value means the vertex lies
upstream of $z = 0$.

The $z$-coordinate of the PoCA point on track 1 is:

\begin{equation}
    z_{\text{PoCA},1} = o_z + s_c \cdot d_z
\end{equation}

\noindent
where $o_z = 0$ is the $z$-component of the reference point $\vec{o}$ (the mathematical intercept anchor of the fit, not the decay vertex location), and $d_z$ is the $z$-component of the unit direction vector $\vec{d}$. For a near-forward track $d_z \approx 1$, so $s_c \approx z_{\text{PoCA},1}$ numerically.
The corresponding transverse position is:

\begin{equation}
    x_{\text{PoCA},1} = o_x + s_c \cdot d_x, \qquad
    y_{\text{PoCA},1} = o_y + s_c \cdot d_y
\end{equation}

\noindent
where $o_x = b_x$ and $o_y = b_y$ are the fitted intercepts, and $d_x$, $d_y$ are the $x$ and $y$ components of the unit direction vector.

% =============================================================================
\subsection{Selecting a Clean Event for Manual Verification}
% =============================================================================

Enable \texttt{DIAG\_RECO\_MODE = true} in \texttt{KLong\_save\_vectors.C} and
scan the printed output for an event satisfying all of the following criteria:

\begin{center}
\begin{tabular}{lll}
    \toprule
    \textbf{Diagnostic} & \textbf{Item in output} & \textbf{Target} \\
    \midrule
    PoCA separation $\Delta$          & \texttt{[1]} & $< 1$ cm \\
    Back-propagated decay time spread & \texttt{[2]} & $|\text{diff}| < 0.1$ ns \\
    Vertex $z$ residual $\delta z$    & \texttt{[3]} & $|\delta z| < 5$ cm \\
    Track slope errors                & \texttt{[4]} & $|\Delta a_x|, |\Delta a_y| < 0.01$ \\
    Tracker coverage                  & hit counts   & T1, T2, T3, T4 all $> 0$ \\
    \bottomrule
\end{tabular}
\end{center}

From a clean event, record the following from the \texttt{DIAG} printout:
\begin{itemize}
    \item All raw hit positions per detector plane (from the \texttt{DEBUG} block).
    \item \texttt{orig} $(b_x, b_y, 0)$ and \texttt{dir} $(d_x, d_y, d_z)$ for each fitted track.
    \item The printed \texttt{reco\_vertex} and \texttt{true\_vertex}.
\end{itemize}

% =============================================================================
\subsection{Manual Calculation Walkthrough}
% =============================================================================

Given the track parameters extracted from a DIAG event for the $\pi^+$ and
$\pi^-$ tracks:

\begin{align}
    \vec{p}_{1} &= (b_x^{\pi^+},\; b_y^{\pi^+},\; 0), &
    \vec{v}_{1} &= (a_x^{\pi^+},\; a_y^{\pi^+},\; 1)^{\text{unit}} \\
    \vec{p}_{2} &= (b_x^{\pi^-},\; b_y^{\pi^-},\; 0), &
    \vec{v}_{2} &= (a_x^{\pi^-},\; a_y^{\pi^-},\; 1)^{\text{unit}}
\end{align}

Proceed step by step:

\begin{enumerate}
    \item Compute $\vec{w}_0 = \vec{p}_1 - \vec{p}_2$.

    \item Since $\vec{v}_{1,2}$ are unit vectors: $a = c = 1$, so
    \begin{equation}
        D = 1 - b^2 = \sin^2\theta,
    \end{equation}
    where $\theta$ is the opening angle between the $\pi^+$ and $\pi^-$ tracks --- the angle between their fitted direction vectors in 3D --- and $b = \cos\theta = \vec{v}_1 \cdot \vec{v}_2$ is its cosine, computed directly from the track fit outputs.

    \item Compute:
    \begin{equation}
        d = \vec{v}_1 \cdot \vec{w}_0, \qquad e = \vec{v}_2 \cdot \vec{w}_0
    \end{equation}

    \item Solve for the PoCA parameters:
    \begin{equation}
        s_c = \frac{be - d}{D}, \qquad t_c = \frac{e - bd}{D}
    \end{equation}

    \item Compute the two closest points:
    \begin{equation}
        \vec{q}_1 = \vec{p}_1 + s_c\,\vec{v}_1, \qquad
        \vec{q}_2 = \vec{p}_2 + t_c\,\vec{v}_2
    \end{equation}

    \item The decay vertex estimate is:
    \begin{equation}
        \vec{V}_\text{reco} = \tfrac{1}{2}\bigl(\vec{q}_1 + \vec{q}_2\bigr)
    \end{equation}

    This should match the \texttt{reco\_vertex} in the DIAG output to three
    decimal places.
\end{enumerate}

% =============================================================================
\subsection{Why \texorpdfstring{$z$}{z} is the Hard Coordinate}
% =============================================================================

The transverse coordinates $x$ and $y$ at the vertex are
\textbf{directly constrained} by the T1 and T2 measurements respectively.
The vertex $z$-position, however, is \textbf{never directly measured} --- it
is determined entirely by the fitted slopes $a_x, a_y$ and the PoCA algebra.

The $z$ error propagates approximately as:

\begin{equation}
    \delta z_\text{vtx} \;\sim\; \frac{\sigma_\text{slope}}{\Delta a^\text{opening}}
\end{equation}

where $\Delta a^\text{opening}$ is the difference in slopes between the
$\pi^+$ and $\pi^-$ tracks and $\sigma_\text{slope}$ is the typical
uncertainty on the fitted slope.

When the opening angle is small (high-momentum kaon, both pions nearly
collinear), $\Delta a^\text{opening}$ is tiny and the $z$ uncertainty diverges.
This is precisely why the vertex summary shows:

\begin{center}
\begin{tabular}{lcc}
    \toprule
    \textbf{Coordinate} & $\mu$ (cm) & $\sigma$ (cm) \\
    \midrule
    $\delta x$ & $\sim -0.05$ & $\sim 0.4$ \\
    $\delta y$ & $\sim 0.00$  & $\sim 0.35$ \\
    $\delta z$ & $\sim +1$--5 & $\sim 7$--12 \\
    \bottomrule
\end{tabular}
\end{center}

\noindent The $z$ resolution is 20--30$\times$ worse than the transverse
resolution, and the positive $\mu_z$ bias (downstream shift) indicates
systematic overestimation of the vertex $z$ position, which propagates
directly into an underestimation of the reconstructed kaon momentum.

% =============================================================================
\subsection{Quick-Start Verification Recipe}
% =============================================================================

The most direct path to a single verified event is:

\begin{enumerate}
    \item In \texttt{KLong\_save\_vectors.C}, set:
    \begin{verbatim}
const bool DEBUG_MODE       = true;
const int  DEBUG_MAX_EVENTS = 1;
const bool DIAG_RECO_MODE   = true;
const int  DIAG_MAX_EVENTS  = 1;
    \end{verbatim}

    \item Run on one simulation file. The first successfully reconstructed
    event will print the full hit list, fitted track parameters, and the
    5-item diagnostic checklist.

    \item Extract \texttt{orig} and \texttt{dir} for each pion track from
    the output.

    \item Work through the 6-step PoCA calculation in Section~5.

    \item Compare your computed $\vec{V}_\text{reco}$ to the printed
    \texttt{reco\_vertex} --- they should agree to $\lesssim 10^{-3}$ cm.

    \item Compare $\vec{V}_\text{reco}$ to the printed \texttt{true\_vertex}
    (from Geant4 MC truth) to quantify the reconstruction resolution for
    that event.
\end{enumerate}

% Required packages in parent preamble:
%   \usepackage{amsmath, amssymb, booktabs, array}
% =============================================================================

% =============================================================================
\subsection{Track Parameterisation}
% =============================================================================

Each pion track is a straight line in 3D, fitted by two independent
linear regressions: one in the $xz$-plane and one in the $yz$-plane.

\begin{equation}
    x(z) = b_x + a_x z, \qquad y(z) = b_y + a_y z
\end{equation}

\noindent
Here $a_x = dx/dz$ and $a_y = dy/dz$ are the fitted track slopes (the rate of transverse displacement per unit distance along the beam axis), and $b_x$, $b_y$ are the track intercepts --- the transverse positions of the track at $z = 0$.

The unit direction vector and reference point (evaluated at $z = 0$, used purely as a mathematical anchor for the fit --- not a physical location) are:

\begin{equation}
    \vec{d} = \bigl(a_x,\; a_y,\; 1\bigr)^{\!\text{unit}}, \qquad
    \vec{o} = (b_x,\; b_y,\; 0)
\end{equation}

Any point on the track is then $\vec{r}(s) = \vec{o} + s\,\vec{d}$, where $s$ is a scalar arc-length parameter: stepping $s$ forward or backward moves along the track in units of centimetres (since $\vec{d}$ is a unit vector).

\subsubsection{Weighted Least-Squares for \texorpdfstring{$a_x, b_x$}{ax, bx}}

To compute the fitted slopes and intercepts by hand from the hit positions,
use weighted least squares. For $N$ hit points $(z_i, x_i)$ --- where $N$ is the total number of hits contributing to the fit, $z_i$ is the beam-axis position of hit $i$, and $x_i$ is the measured transverse coordinate --- with weights $w_i = 1/\sigma_i^2$, where $\sigma_i$ is the positional resolution of the detector element that produced hit $i$ (values given in the table below):

\begin{equation}
    a_x = \frac{\langle zx \rangle - \langle z \rangle\langle x \rangle}{\langle z^2 \rangle - \langle z \rangle^2},
    \qquad
    b_x = \langle x \rangle - a_x\,\langle z \rangle,
    \qquad \text{where} \quad
    \langle f \rangle = \frac{\sum_i w_i f_i}{\sum_i w_i}
\end{equation}

The same procedure gives $a_y, b_y$ from the $(z_i, y_i)$ points.

\subsubsection{Detector Measurement Uncertainties}

The per-hit uncertainties $\sigma_i$ used as weights in the track fit are:

\begin{center}
\begin{tabular}{llll}
    \toprule
    \textbf{Detector} & \textbf{Device IDs} & \textbf{Axis measured} & \textbf{$\sigma$ (cm)} \\
    \midrule
    T1 (straws along $Y$)        & 1--488      & $X$ only            & 0.40 \\
    T2 (straws along $X$)        & 489--976    & $Y$ only            & 0.40 \\
    T3 (stereo $+45°$)           & 977--1464   & $X$ and $Y$         & $0.40\sqrt{2} \approx 0.566$ \\
    T4 (stereo $-45°$)           & 1465--1952  & $X$ and $Y$         & $0.40\sqrt{2} \approx 0.566$ \\
    FRI Wall 1 (inner strips)    & 2001--2010, 2017--2026 & $X$ only & 1.50 \\
    FRI Wall 1 (outer strips)    & 2011--2016  & $X$ only            & 3.00 \\
    FRI Wall 2 (inner strips)    & 2027--2036, 2043--2052 & $Y$ only & 1.50 \\
    FRI Wall 2 (outer strips)    & 2037--2042  & $Y$ only            & 3.00 \\
    TOF wall bars                & 2053--2070  & \textit{excluded from fit} & --- \\
    \bottomrule
\end{tabular}
\end{center}

\noindent
The TOF bars are excluded from the track fit because their large positional
uncertainties ($\pm 6$ cm in $X$, $\pm 10$ cm in $Y$) at the $\sim\!460$ cm
lever arm would bias the fitted slope and degrade backward-extrapolated vertex
resolution. TOF hits are used only for the velocity measurement.

% =============================================================================
\subsection{Point of Closest Approach (PoCA) --- Full Derivation}
% =============================================================================

Given two fitted straight-line tracks:

\begin{equation}
    \vec{r}_1(s) = \vec{p}_1 + s\,\vec{v}_1, \qquad
    \vec{r}_2(t) = \vec{p}_2 + t\,\vec{v}_2,
\end{equation}

\noindent
where $\vec{p}_{1,2}$ are the mathematical reference points on each track, defined as the position where the fitted line passes through the $z = 0$ plane (this is not the decay vertex or the kaon production point --- it is simply the intercept of the extrapolated track line, used as a convenient algebraic anchor), $\vec{v}_{1,2}$ are the unit direction vectors for the $\pi^+$ and $\pi^-$ tracks respectively (corresponding to $\vec{d}$ above), and $s$, $t$ are free scalar arc-length parameters that index positions along each track. The goal is to find the values of $s$ and $t$ that minimise the distance $\bigl|\vec{r}_1(s) - \vec{r}_2(t)\bigr|$.

\subsubsection{Setup}

Define:

\begin{equation}
    \vec{w}_0 = \vec{p}_1 - \vec{p}_2
\end{equation}

\noindent
The vector $\vec{w}_0$ is the displacement from the $\pi^-$ reference point to the $\pi^+$ reference point; it encodes the initial separation between the two track lines at $z = 0$.

Define also the five scalar dot products:

\begin{align}
    a &= \vec{v}_1 \cdot \vec{v}_1, &
    b &= \vec{v}_1 \cdot \vec{v}_2, &
    c &= \vec{v}_2 \cdot \vec{v}_2, \\[4pt]
    d &= \vec{v}_1 \cdot \vec{w}_0, &
    e &= \vec{v}_2 \cdot \vec{w}_0. & &
\end{align}

\noindent
Since $\vec{v}_{1,2}$ are unit vectors, $a = c = 1$ always. The quantity $b = \cos\theta$ is the cosine of the opening angle $\theta$ between the two pion tracks. The scalars $d$ and $e$ project the inter-track separation vector $\vec{w}_0$ onto the $\pi^+$ and $\pi^-$ direction vectors respectively; they determine how the gap between the reference points must be closed by stepping along each track.

\subsubsection{Minimisation}

Setting $\partial/\partial s$ and $\partial/\partial t$ of
$|\vec{r}_1(s) - \vec{r}_2(t)|^2$ to zero gives the $2 \times 2$ linear system:

\begin{equation}
    \begin{pmatrix} a & -b \\ b & -c \end{pmatrix}
    \begin{pmatrix} s_c \\ t_c \end{pmatrix}
    = \begin{pmatrix} -d \\ -e \end{pmatrix}
\end{equation}

Solving (Cramer's rule), with denominator $D = ac - b^2 = 1 - \cos^2\theta = \sin^2\theta$:

\begin{equation}
    \boxed{s_c = \frac{be - cd}{D}, \qquad t_c = \frac{ae - bd}{D}}
\end{equation}

\noindent
The denominator $D = \sin^2\theta$ is a measure of how non-parallel the two tracks are; it reaches its maximum of 1 for perpendicular tracks and vanishes as the opening angle $\theta \to 0$. The solved parameters $s_c$ and $t_c$ are the arc-length values at which each track is closest to the other --- stepping $s_c$ along track 1 and $t_c$ along track 2 brings both tracks to their point of minimum mutual separation. If $|D| < 10^{-14}$ the lines are parallel and no unique solution exists.

\subsubsection{Closest Points and Decay Vertex}

The closest point on each track is:

\begin{equation}
    \vec{q}_1 = \vec{p}_1 + s_c\,\vec{v}_1, \qquad
    \vec{q}_2 = \vec{p}_2 + t_c\,\vec{v}_2
\end{equation}

\noindent
The points $\vec{q}_1$ and $\vec{q}_2$ are the feet of the common perpendicular between the two track lines --- the positions on the $\pi^+$ and $\pi^-$ tracks respectively that are nearest to each other. In a perfect detector with infinitely precise tracking, both pions would originate from the same point and $\vec{q}_1 = \vec{q}_2$ exactly.

The reconstructed decay vertex is the midpoint of the shortest connecting
segment:

\begin{equation}
    \boxed{\vec{V}_\text{reco} = \tfrac{1}{2}\bigl(\vec{q}_1 + \vec{q}_2\bigr)}
\end{equation}

The \emph{PoCA separation} (a quality indicator) is:

\begin{equation}
    \Delta = |\vec{q}_1 - \vec{q}_2|
\end{equation}

For a clean decay event occurring at a well-defined vertex, $\Delta$ should be
close to zero. Large $\Delta$ indicates nearly-parallel tracks, a track mismatch,
or a poor fit.

% =============================================================================
\subsection{Physical Interpretation of \texorpdfstring{$s_c$}{sc} and \texorpdfstring{$t_c$}{tc}}
% =============================================================================

Because $\vec{v}_{1,2}$ are unit vectors ($|\vec{v}| = 1$), the parameters
$s_c$ and $t_c$ are the \textbf{arc-lengths along each track} from the reference
point $\vec{o}$ to the PoCA point. A negative value means the vertex lies
upstream of $z = 0$.

The $z$-coordinate of the PoCA point on track 1 is:

\begin{equation}
    z_{\text{PoCA},1} = o_z + s_c \cdot d_z
\end{equation}

\noindent
where $o_z = 0$ is the $z$-component of the reference point $\vec{o}$ (the mathematical intercept anchor of the fit, not the decay vertex location), and $d_z$ is the $z$-component of the unit direction vector $\vec{d}$. For a near-forward track $d_z \approx 1$, so $s_c \approx z_{\text{PoCA},1}$ numerically.
The corresponding transverse position is:

\begin{equation}
    x_{\text{PoCA},1} = o_x + s_c \cdot d_x, \qquad
    y_{\text{PoCA},1} = o_y + s_c \cdot d_y
\end{equation}

\noindent
where $o_x = b_x$ and $o_y = b_y$ are the fitted intercepts, and $d_x$, $d_y$ are the $x$ and $y$ components of the unit direction vector.

% =============================================================================
\subsection{Selecting a Clean Event for Manual Verification}
% =============================================================================

Enable \texttt{DIAG\_RECO\_MODE = true} in \texttt{KLong\_save\_vectors.C} and
scan the printed output for an event satisfying all of the following criteria:

\begin{center}
\begin{tabular}{lll}
    \toprule
    \textbf{Diagnostic} & \textbf{Item in output} & \textbf{Target} \\
    \midrule
    PoCA separation $\Delta$          & \texttt{[1]} & $< 1$ cm \\
    Back-propagated decay time spread & \texttt{[2]} & $|\text{diff}| < 0.1$ ns \\
    Vertex $z$ residual $\delta z$    & \texttt{[3]} & $|\delta z| < 5$ cm \\
    Track slope errors                & \texttt{[4]} & $|\Delta a_x|, |\Delta a_y| < 0.01$ \\
    Tracker coverage                  & hit counts   & T1, T2, T3, T4 all $> 0$ \\
    \bottomrule
\end{tabular}
\end{center}

From a clean event, record the following from the \texttt{DIAG} printout:
\begin{itemize}
    \item All raw hit positions per detector plane (from the \texttt{DEBUG} block).
    \item \texttt{orig} $(b_x, b_y, 0)$ and \texttt{dir} $(d_x, d_y, d_z)$ for each fitted track.
    \item The printed \texttt{reco\_vertex} and \texttt{true\_vertex}.
\end{itemize}

% =============================================================================
\subsection{Manual Calculation Walkthrough}
% =============================================================================

Given the track parameters extracted from a DIAG event for the $\pi^+$ and
$\pi^-$ tracks:

\begin{align}
    \vec{p}_{1} &= (b_x^{\pi^+},\; b_y^{\pi^+},\; 0), &
    \vec{v}_{1} &= (a_x^{\pi^+},\; a_y^{\pi^+},\; 1)^{\text{unit}} \\
    \vec{p}_{2} &= (b_x^{\pi^-},\; b_y^{\pi^-},\; 0), &
    \vec{v}_{2} &= (a_x^{\pi^-},\; a_y^{\pi^-},\; 1)^{\text{unit}}
\end{align}

Proceed step by step:

\begin{enumerate}
    \item Compute $\vec{w}_0 = \vec{p}_1 - \vec{p}_2$.

    \item Since $\vec{v}_{1,2}$ are unit vectors: $a = c = 1$, so
    \begin{equation}
        D = 1 - b^2 = \sin^2\theta,
    \end{equation}
    where $\theta$ is the opening angle between the $\pi^+$ and $\pi^-$ tracks --- the angle between their fitted direction vectors in 3D --- and $b = \cos\theta = \vec{v}_1 \cdot \vec{v}_2$ is its cosine, computed directly from the track fit outputs.

    \item Compute:
    \begin{equation}
        d = \vec{v}_1 \cdot \vec{w}_0, \qquad e = \vec{v}_2 \cdot \vec{w}_0
    \end{equation}

    \item Solve for the PoCA parameters:
    \begin{equation}
        s_c = \frac{be - d}{D}, \qquad t_c = \frac{e - bd}{D}
    \end{equation}

    \item Compute the two closest points:
    \begin{equation}
        \vec{q}_1 = \vec{p}_1 + s_c\,\vec{v}_1, \qquad
        \vec{q}_2 = \vec{p}_2 + t_c\,\vec{v}_2
    \end{equation}

    \item The decay vertex estimate is:
    \begin{equation}
        \vec{V}_\text{reco} = \tfrac{1}{2}\bigl(\vec{q}_1 + \vec{q}_2\bigr)
    \end{equation}

    This should match the \texttt{reco\_vertex} in the DIAG output to three
    decimal places.
\end{enumerate}

% =============================================================================
\subsection{Why \texorpdfstring{$z$}{z} is the Hard Coordinate}
% =============================================================================

The transverse coordinates $x$ and $y$ at the vertex are
\textbf{directly constrained} by the T1 and T2 measurements respectively.
The vertex $z$-position, however, is \textbf{never directly measured} --- it
is determined entirely by the fitted slopes $a_x, a_y$ and the PoCA algebra.

The $z$ error propagates approximately as:

\begin{equation}
    \delta z_\text{vtx} \;\sim\; \frac{\sigma_\text{slope}}{\Delta a^\text{opening}}
\end{equation}

where $\Delta a^\text{opening}$ is the difference in slopes between the
$\pi^+$ and $\pi^-$ tracks and $\sigma_\text{slope}$ is the typical
uncertainty on the fitted slope.

When the opening angle is small (high-momentum kaon, both pions nearly
collinear), $\Delta a^\text{opening}$ is tiny and the $z$ uncertainty diverges.
This is precisely why the vertex summary shows:

\begin{center}
\begin{tabular}{lcc}
    \toprule
    \textbf{Coordinate} & $\mu$ (cm) & $\sigma$ (cm) \\
    \midrule
    $\delta x$ & $\sim -0.05$ & $\sim 0.4$ \\
    $\delta y$ & $\sim 0.00$  & $\sim 0.35$ \\
    $\delta z$ & $\sim +1$--5 & $\sim 7$--12 \\
    \bottomrule
\end{tabular}
\end{center}

\noindent The $z$ resolution is 20--30$\times$ worse than the transverse
resolution, and the positive $\mu_z$ bias (downstream shift) indicates
systematic overestimation of the vertex $z$ position, which propagates
directly into an underestimation of the reconstructed kaon momentum.

% =============================================================================
\subsection{Quick-Start Verification Recipe}
% =============================================================================

The most direct path to a single verified event is:

\begin{enumerate}
    \item In \texttt{KLong\_save\_vectors.C}, set:
    \begin{verbatim}
const bool DEBUG_MODE       = true;
const int  DEBUG_MAX_EVENTS = 1;
const bool DIAG_RECO_MODE   = true;
const int  DIAG_MAX_EVENTS  = 1;
    \end{verbatim}

    \item Run on one simulation file. The first successfully reconstructed
    event will print the full hit list, fitted track parameters, and the
    5-item diagnostic checklist.

    \item Extract \texttt{orig} and \texttt{dir} for each pion track from
    the output.

    \item Work through the 6-step PoCA calculation in Section~5.

    \item Compare your computed $\vec{V}_\text{reco}$ to the printed
    \texttt{reco\_vertex} --- they should agree to $\lesssim 10^{-3}$ cm.

    \item Compare $\vec{V}_\text{reco}$ to the printed \texttt{true\_vertex}
    (from Geant4 MC truth) to quantify the reconstruction resolution for
    that event.
\end{enumerate}

% Required packages in parent preamble:
%   \usepackage{amsmath, amssymb, booktabs, array}
% =============================================================================

% =============================================================================
\subsection{Track Parameterisation}
% =============================================================================

Each pion track is a straight line in 3D, fitted by two independent
linear regressions: one in the $xz$-plane and one in the $yz$-plane.

\begin{equation}
    x(z) = b_x + a_x z, \qquad y(z) = b_y + a_y z
\end{equation}

\noindent
Here $a_x = dx/dz$ and $a_y = dy/dz$ are the fitted track slopes (the rate of transverse displacement per unit distance along the beam axis), and $b_x$, $b_y$ are the track intercepts --- the transverse positions of the track at $z = 0$.

The unit direction vector and reference point (evaluated at $z = 0$, used purely as a mathematical anchor for the fit --- not a physical location) are:

\begin{equation}
    \vec{d} = \bigl(a_x,\; a_y,\; 1\bigr)^{\!\text{unit}}, \qquad
    \vec{o} = (b_x,\; b_y,\; 0)
\end{equation}

Any point on the track is then $\vec{r}(s) = \vec{o} + s\,\vec{d}$, where $s$ is a scalar arc-length parameter: stepping $s$ forward or backward moves along the track in units of centimetres (since $\vec{d}$ is a unit vector).

\subsubsection{Weighted Least-Squares for \texorpdfstring{$a_x, b_x$}{ax, bx}}

To compute the fitted slopes and intercepts by hand from the hit positions,
use weighted least squares. For $N$ hit points $(z_i, x_i)$ --- where $N$ is the total number of hits contributing to the fit, $z_i$ is the beam-axis position of hit $i$, and $x_i$ is the measured transverse coordinate --- with weights $w_i = 1/\sigma_i^2$, where $\sigma_i$ is the positional resolution of the detector element that produced hit $i$ (values given in the table below):

\begin{equation}
    a_x = \frac{\langle zx \rangle - \langle z \rangle\langle x \rangle}{\langle z^2 \rangle - \langle z \rangle^2},
    \qquad
    b_x = \langle x \rangle - a_x\,\langle z \rangle,
    \qquad \text{where} \quad
    \langle f \rangle = \frac{\sum_i w_i f_i}{\sum_i w_i}
\end{equation}

The same procedure gives $a_y, b_y$ from the $(z_i, y_i)$ points.

\subsubsection{Detector Measurement Uncertainties}

The per-hit uncertainties $\sigma_i$ used as weights in the track fit are:

\begin{center}
\begin{tabular}{llll}
    \toprule
    \textbf{Detector} & \textbf{Device IDs} & \textbf{Axis measured} & \textbf{$\sigma$ (cm)} \\
    \midrule
    T1 (straws along $Y$)        & 1--488      & $X$ only            & 0.40 \\
    T2 (straws along $X$)        & 489--976    & $Y$ only            & 0.40 \\
    T3 (stereo $+45°$)           & 977--1464   & $X$ and $Y$         & $0.40\sqrt{2} \approx 0.566$ \\
    T4 (stereo $-45°$)           & 1465--1952  & $X$ and $Y$         & $0.40\sqrt{2} \approx 0.566$ \\
    FRI Wall 1 (inner strips)    & 2001--2010, 2017--2026 & $X$ only & 1.50 \\
    FRI Wall 1 (outer strips)    & 2011--2016  & $X$ only            & 3.00 \\
    FRI Wall 2 (inner strips)    & 2027--2036, 2043--2052 & $Y$ only & 1.50 \\
    FRI Wall 2 (outer strips)    & 2037--2042  & $Y$ only            & 3.00 \\
    TOF wall bars                & 2053--2070  & \textit{excluded from fit} & --- \\
    \bottomrule
\end{tabular}
\end{center}

\noindent
The TOF bars are excluded from the track fit because their large positional
uncertainties ($\pm 6$ cm in $X$, $\pm 10$ cm in $Y$) at the $\sim\!460$ cm
lever arm would bias the fitted slope and degrade backward-extrapolated vertex
resolution. TOF hits are used only for the velocity measurement.

% =============================================================================
\subsection{Point of Closest Approach (PoCA) --- Full Derivation}
% =============================================================================

Given two fitted straight-line tracks:

\begin{equation}
    \vec{r}_1(s) = \vec{p}_1 + s\,\vec{v}_1, \qquad
    \vec{r}_2(t) = \vec{p}_2 + t\,\vec{v}_2,
\end{equation}

\noindent
where $\vec{p}_{1,2}$ are the mathematical reference points on each track, defined as the position where the fitted line passes through the $z = 0$ plane (this is not the decay vertex or the kaon production point --- it is simply the intercept of the extrapolated track line, used as a convenient algebraic anchor), $\vec{v}_{1,2}$ are the unit direction vectors for the $\pi^+$ and $\pi^-$ tracks respectively (corresponding to $\vec{d}$ above), and $s$, $t$ are free scalar arc-length parameters that index positions along each track. The goal is to find the values of $s$ and $t$ that minimise the distance $\bigl|\vec{r}_1(s) - \vec{r}_2(t)\bigr|$.

\subsubsection{Setup}

Define:

\begin{equation}
    \vec{w}_0 = \vec{p}_1 - \vec{p}_2
\end{equation}

\noindent
The vector $\vec{w}_0$ is the displacement from the $\pi^-$ reference point to the $\pi^+$ reference point; it encodes the initial separation between the two track lines at $z = 0$.

Define also the five scalar dot products:

\begin{align}
    a &= \vec{v}_1 \cdot \vec{v}_1, &
    b &= \vec{v}_1 \cdot \vec{v}_2, &
    c &= \vec{v}_2 \cdot \vec{v}_2, \\[4pt]
    d &= \vec{v}_1 \cdot \vec{w}_0, &
    e &= \vec{v}_2 \cdot \vec{w}_0. & &
\end{align}

\noindent
Since $\vec{v}_{1,2}$ are unit vectors, $a = c = 1$ always. The quantity $b = \cos\theta$ is the cosine of the opening angle $\theta$ between the two pion tracks. The scalars $d$ and $e$ project the inter-track separation vector $\vec{w}_0$ onto the $\pi^+$ and $\pi^-$ direction vectors respectively; they determine how the gap between the reference points must be closed by stepping along each track.

\subsubsection{Minimisation}

Setting $\partial/\partial s$ and $\partial/\partial t$ of
$|\vec{r}_1(s) - \vec{r}_2(t)|^2$ to zero gives the $2 \times 2$ linear system:

\begin{equation}
    \begin{pmatrix} a & -b \\ b & -c \end{pmatrix}
    \begin{pmatrix} s_c \\ t_c \end{pmatrix}
    = \begin{pmatrix} -d \\ -e \end{pmatrix}
\end{equation}

Solving (Cramer's rule), with denominator $D = ac - b^2 = 1 - \cos^2\theta = \sin^2\theta$:

\begin{equation}
    \boxed{s_c = \frac{be - cd}{D}, \qquad t_c = \frac{ae - bd}{D}}
\end{equation}

\noindent
The denominator $D = \sin^2\theta$ is a measure of how non-parallel the two tracks are; it reaches its maximum of 1 for perpendicular tracks and vanishes as the opening angle $\theta \to 0$. The solved parameters $s_c$ and $t_c$ are the arc-length values at which each track is closest to the other --- stepping $s_c$ along track 1 and $t_c$ along track 2 brings both tracks to their point of minimum mutual separation. If $|D| < 10^{-14}$ the lines are parallel and no unique solution exists.

\subsubsection{Closest Points and Decay Vertex}

The closest point on each track is:

\begin{equation}
    \vec{q}_1 = \vec{p}_1 + s_c\,\vec{v}_1, \qquad
    \vec{q}_2 = \vec{p}_2 + t_c\,\vec{v}_2
\end{equation}

\noindent
The points $\vec{q}_1$ and $\vec{q}_2$ are the feet of the common perpendicular between the two track lines --- the positions on the $\pi^+$ and $\pi^-$ tracks respectively that are nearest to each other. In a perfect detector with infinitely precise tracking, both pions would originate from the same point and $\vec{q}_1 = \vec{q}_2$ exactly.

The reconstructed decay vertex is the midpoint of the shortest connecting
segment:

\begin{equation}
    \boxed{\vec{V}_\text{reco} = \tfrac{1}{2}\bigl(\vec{q}_1 + \vec{q}_2\bigr)}
\end{equation}

The \emph{PoCA separation} (a quality indicator) is:

\begin{equation}
    \Delta = |\vec{q}_1 - \vec{q}_2|
\end{equation}

For a clean decay event occurring at a well-defined vertex, $\Delta$ should be
close to zero. Large $\Delta$ indicates nearly-parallel tracks, a track mismatch,
or a poor fit.

% =============================================================================
\subsection{Physical Interpretation of \texorpdfstring{$s_c$}{sc} and \texorpdfstring{$t_c$}{tc}}
% =============================================================================

Because $\vec{v}_{1,2}$ are unit vectors ($|\vec{v}| = 1$), the parameters
$s_c$ and $t_c$ are the \textbf{arc-lengths along each track} from the reference
point $\vec{o}$ to the PoCA point. A negative value means the vertex lies
upstream of $z = 0$.

The $z$-coordinate of the PoCA point on track 1 is:

\begin{equation}
    z_{\text{PoCA},1} = o_z + s_c \cdot d_z
\end{equation}

\noindent
where $o_z = 0$ is the $z$-component of the reference point $\vec{o}$ (the mathematical intercept anchor of the fit, not the decay vertex location), and $d_z$ is the $z$-component of the unit direction vector $\vec{d}$. For a near-forward track $d_z \approx 1$, so $s_c \approx z_{\text{PoCA},1}$ numerically.
The corresponding transverse position is:

\begin{equation}
    x_{\text{PoCA},1} = o_x + s_c \cdot d_x, \qquad
    y_{\text{PoCA},1} = o_y + s_c \cdot d_y
\end{equation}

\noindent
where $o_x = b_x$ and $o_y = b_y$ are the fitted intercepts, and $d_x$, $d_y$ are the $x$ and $y$ components of the unit direction vector.

% =============================================================================
\subsection{Selecting a Clean Event for Manual Verification}
% =============================================================================

Enable \texttt{DIAG\_RECO\_MODE = true} in \texttt{KLong\_save\_vectors.C} and
scan the printed output for an event satisfying all of the following criteria:

\begin{center}
\begin{tabular}{lll}
    \toprule
    \textbf{Diagnostic} & \textbf{Item in output} & \textbf{Target} \\
    \midrule
    PoCA separation $\Delta$          & \texttt{[1]} & $< 1$ cm \\
    Back-propagated decay time spread & \texttt{[2]} & $|\text{diff}| < 0.1$ ns \\
    Vertex $z$ residual $\delta z$    & \texttt{[3]} & $|\delta z| < 5$ cm \\
    Track slope errors                & \texttt{[4]} & $|\Delta a_x|, |\Delta a_y| < 0.01$ \\
    Tracker coverage                  & hit counts   & T1, T2, T3, T4 all $> 0$ \\
    \bottomrule
\end{tabular}
\end{center}

From a clean event, record the following from the \texttt{DIAG} printout:
\begin{itemize}
    \item All raw hit positions per detector plane (from the \texttt{DEBUG} block).
    \item \texttt{orig} $(b_x, b_y, 0)$ and \texttt{dir} $(d_x, d_y, d_z)$ for each fitted track.
    \item The printed \texttt{reco\_vertex} and \texttt{true\_vertex}.
\end{itemize}

% =============================================================================
\subsection{Manual Calculation Walkthrough}
% =============================================================================

Given the track parameters extracted from a DIAG event for the $\pi^+$ and
$\pi^-$ tracks:

\begin{align}
    \vec{p}_{1} &= (b_x^{\pi^+},\; b_y^{\pi^+},\; 0), &
    \vec{v}_{1} &= (a_x^{\pi^+},\; a_y^{\pi^+},\; 1)^{\text{unit}} \\
    \vec{p}_{2} &= (b_x^{\pi^-},\; b_y^{\pi^-},\; 0), &
    \vec{v}_{2} &= (a_x^{\pi^-},\; a_y^{\pi^-},\; 1)^{\text{unit}}
\end{align}

Proceed step by step:

\begin{enumerate}
    \item Compute $\vec{w}_0 = \vec{p}_1 - \vec{p}_2$.

    \item Since $\vec{v}_{1,2}$ are unit vectors: $a = c = 1$, so
    \begin{equation}
        D = 1 - b^2 = \sin^2\theta,
    \end{equation}
    where $\theta$ is the opening angle between the $\pi^+$ and $\pi^-$ tracks --- the angle between their fitted direction vectors in 3D --- and $b = \cos\theta = \vec{v}_1 \cdot \vec{v}_2$ is its cosine, computed directly from the track fit outputs.

    \item Compute:
    \begin{equation}
        d = \vec{v}_1 \cdot \vec{w}_0, \qquad e = \vec{v}_2 \cdot \vec{w}_0
    \end{equation}

    \item Solve for the PoCA parameters:
    \begin{equation}
        s_c = \frac{be - d}{D}, \qquad t_c = \frac{e - bd}{D}
    \end{equation}

    \item Compute the two closest points:
    \begin{equation}
        \vec{q}_1 = \vec{p}_1 + s_c\,\vec{v}_1, \qquad
        \vec{q}_2 = \vec{p}_2 + t_c\,\vec{v}_2
    \end{equation}

    \item The decay vertex estimate is:
    \begin{equation}
        \vec{V}_\text{reco} = \tfrac{1}{2}\bigl(\vec{q}_1 + \vec{q}_2\bigr)
    \end{equation}

    This should match the \texttt{reco\_vertex} in the DIAG output to three
    decimal places.
\end{enumerate}

% =============================================================================
\subsection{Why \texorpdfstring{$z$}{z} is the Hard Coordinate}
% =============================================================================

The transverse coordinates $x$ and $y$ at the vertex are
\textbf{directly constrained} by the T1 and T2 measurements respectively.
The vertex $z$-position, however, is \textbf{never directly measured} --- it
is determined entirely by the fitted slopes $a_x, a_y$ and the PoCA algebra.

The $z$ error propagates approximately as:

\begin{equation}
    \delta z_\text{vtx} \;\sim\; \frac{\sigma_\text{slope}}{\Delta a^\text{opening}}
\end{equation}

where $\Delta a^\text{opening}$ is the difference in slopes between the
$\pi^+$ and $\pi^-$ tracks and $\sigma_\text{slope}$ is the typical
uncertainty on the fitted slope.

When the opening angle is small (high-momentum kaon, both pions nearly
collinear), $\Delta a^\text{opening}$ is tiny and the $z$ uncertainty diverges.
This is precisely why the vertex summary shows:

\begin{center}
\begin{tabular}{lcc}
    \toprule
    \textbf{Coordinate} & $\mu$ (cm) & $\sigma$ (cm) \\
    \midrule
    $\delta x$ & $\sim -0.05$ & $\sim 0.4$ \\
    $\delta y$ & $\sim 0.00$  & $\sim 0.35$ \\
    $\delta z$ & $\sim +1$--5 & $\sim 7$--12 \\
    \bottomrule
\end{tabular}
\end{center}

\noindent The $z$ resolution is 20--30$\times$ worse than the transverse
resolution, and the positive $\mu_z$ bias (downstream shift) indicates
systematic overestimation of the vertex $z$ position, which propagates
directly into an underestimation of the reconstructed kaon momentum.

% =============================================================================
\subsection{Quick-Start Verification Recipe}
% =============================================================================

The most direct path to a single verified event is:

\begin{enumerate}
    \item In \texttt{KLong\_save\_vectors.C}, set:
    \begin{verbatim}
const bool DEBUG_MODE       = true;
const int  DEBUG_MAX_EVENTS = 1;
const bool DIAG_RECO_MODE   = true;
const int  DIAG_MAX_EVENTS  = 1;
    \end{verbatim}

    \item Run on one simulation file. The first successfully reconstructed
    event will print the full hit list, fitted track parameters, and the
    5-item diagnostic checklist.

    \item Extract \texttt{orig} and \texttt{dir} for each pion track from
    the output.

    \item Work through the 6-step PoCA calculation in Section~5.

    \item Compare your computed $\vec{V}_\text{reco}$ to the printed
    \texttt{reco\_vertex} --- they should agree to $\lesssim 10^{-3}$ cm.

    \item Compare $\vec{V}_\text{reco}$ to the printed \texttt{true\_vertex}
    (from Geant4 MC truth) to quantify the reconstruction resolution for
    that event.
\end{enumerate}
